\section{Blind Damage: Theory against Data}
\label{sec:blind}

\subsection{Connectivity}
\label{sec:blind-connectivity}

Under random edge removal the giant component barely shrinks until most edges
are gone (Figure~\ref{fig:summary}a). The ten-order means of $S$ are 0.982 at
$f=0.5$, 0.815 at $f=0.9$ and 0.112 at $f=0.99$. Their sample standard
deviations are 0.001--0.016. Random node removal gives $S=0.493$ at $f=0.5$,
which is 98.5\% of the surviving nodes, so the removed nodes account for almost
all of the loss. The generating-function predictions,
Eqs.~\eqref{eq:gf-bond}--\eqref{eq:gf-site}, stay within 0.0011 (edges) and
0.0033 (nodes) of these means for $f\le0.9$. As Section~\ref{sec:theory-gf}
showed, this agreement mainly confirms the local limit, Eq.~\eqref{eq:s-local}:
a node drops out when it loses its last edge.

Two checks bound what this says about PrimeKG. First, the solver itself was
tested on ten random graphs built by matching the stubs of the same degree
sequence, where it agrees with simulation to within 0.0012 for edge removal and
node removal, and for the degree attack up to $f=0.5$. The union-find simulation was also compared
with direct NetworkX calculations at every damage state of 80 small graphs.
Second, the filtered full graph is far less robust: at $f=0.5$ and $f=0.9$ its
giant component holds 0.858 and 0.554 of nodes, against 0.982 and 0.815 for the
sample. Its mean degree is lower (27.6 against 47.2), so more nodes have few
edges, which is Eq.~\eqref{eq:s-local} again. The connectivity numbers in this
report therefore describe the dense BFS sample. The full graph is less robust.

\subsection{Two-hop retrieval}
\label{sec:blind-r2}

Figure~\ref{fig:r2-long-run} compares Eq.~\eqref{eq:r2} with the mean over 5,000
independent damage orders for each query group. At $f=0.5$ the single-path group
retains 0.250 of its queries (formula $0.5^2=0.250$), and the multi-path group
0.576 (formula 0.577). In a separate run of 5,000 orders for the whole
benchmark, the largest difference from Eq.~\eqref{eq:r2} at any damage level is
0.0008 for edge removal and 0.0010 for node removal, at most 2.4 Monte Carlo
standard errors. As
Section~\ref{sec:theory-exact} explains, this confirms the implementation and
the independence of routes. Its practical value is that the two-hop
robustness of any set of queries can be computed from its redundancy
distribution alone, without simulation.

We use this to look beyond our benchmark, which was balanced to contain equal
numbers of single-path and multi-path queries. Among all 4,439 drug--disease
pairs with a shared protein in the standard graph, 71.5\% have $m=1$, and
Eq.~\eqref{eq:r2} gives $R_2=0.35$ at $f=0.5$. Among the 406,600 such pairs in
the full filtered graph, 81\% have $m=1$ and $R_2=0.30$. For a typical
drug--disease association in PrimeKG, then, fewer than a third keep a two-hop
link after half of the edges are lost. The balanced benchmark overstates the
two-hop robustness of typical pairs.

\begin{figure}[!htbp]
  \centering
  \includegraphics[width=0.92\linewidth]{act2/primekg_5000_bfs_v1/fig_R2_long_run.pdf}
  \caption{Two-hop retrieval under random edge removal by query group. Solid
  lines are means over 5,000 independent damage orders. Shading shows the spread
  of a single damage state (one standard deviation). Dashed lines are
  Eq.~\eqref{eq:r2}.}
  \label{fig:r2-long-run}
\end{figure}

Redundancy matters a great deal at moderate damage and little at heavy damage.
At $f=0.3$ a second shared protein raises survival from 0.49 to 0.74, but at
$f=0.7$ only from 0.09 to 0.17, because each extra route adds the same factor
$1-p^2$, which approaches 1. Blind node removal gives $R_2=0.168$ at $\phi=0.5$,
matching Eq.~\eqref{eq:r2-node}. Here the endpoints themselves are lost with
probability $1-\phi^2=0.75$.

\subsection{The horizon and the bottleneck}
\label{sec:blind-horizon}

At $f=0.5$ the long-run means are $S=0.982$, $R_\infty=0.862$, $R_4=0.861$,
$R_3=0.792$ and $R_2=0.414$. The 57-point gap between connectivity and two-hop
retrieval therefore has two parts. About 12 points come from the step between
counting nodes and counting queries: a query needs both of its endpoints
attached, and its drug often has few links. The remaining 45 points come from
the hop limit. Most of the gap is a statement about bounded paths, the regime of
limited-path percolation, rather than about fragmentation.

Each additional hop moves the damage at which the limit starts to matter. We
take this point to be the first $f$ at which $R_L$ falls 5 points below
$R_\infty$. It is $f\approx0.10$, 0.50 and 0.85 for $L=2$, 3 and 4, from the 600
independent orders.

For unlimited hops, Eq.~\eqref{eq:rinf} agrees with the long-run means to within
0.009 at every damage level from $f=0.05$ to $0.99$ (Figure~\ref{fig:summary}a).
The simpler min-cut form, Eq.~\eqref{eq:mincut}, agrees to within 0.005 for
$f\le0.8$. Beyond that it overestimates success, by 0.021 at $f=0.9$, because
the PPI giant component begins to shrink. Even so, the PPI layer outlasts the
query endpoints. Its giant component still holds 75\% of proteins at $f=0.9$
(generating-function value 0.751), and its percolation threshold, Eq.~\eqref{eq:pc}, lies beyond
99\% edge loss. Up to $f=0.9$, more than 90\% of the queries that fail at
$L=\infty$ have lost every link of their drug while the disease side is still
attached. Longer horizons therefore move the bottleneck from shared proteins to
drug target links.

This bottleneck depends on the drug degree distribution. In the sample, 69\% of
drugs have one target. In the full filtered graph the figure is 53\%, so drug
links are a weak point of PrimeKG as a whole, but the size of the effect here
is specific to the sample.

\subsection{How certain are the numbers?}
\label{sec:blind-uncertainty}

Our figures show wide bands, yet we report agreement within 0.001. These
statements answer different questions, so we separate three measures of
uncertainty (Table~\ref{tab:uncertainty}).

\begin{table}[!htbp]
\centering
\small
\caption{Uncertainty in $R_2$ under random edge removal. ``Independent'' is the
standard deviation expected if the 200 queries failed independently with their
formula probabilities.}
\label{tab:uncertainty}
\begin{tabular}{@{}crrrrrr@{}}
\toprule
& & \multicolumn{2}{c}{SD of one damage state} & & \multicolumn{2}{c}{Mean (± standard error)} \\
\cmidrule(lr){3-4}\cmidrule(lr){6-7}
$f$ & Eq.~\eqref{eq:r2} & Observed & Independent & Variance ratio & 10 orders & 5,000 orders \\
\midrule
0.3 & 0.663 & 0.046 & 0.031 & 2.2 & $0.680 \pm 0.012$ & $0.6628 \pm 0.0006$ \\
0.5 & 0.413 & 0.047 & 0.032 & 2.2 & $0.433 \pm 0.013$ & $0.4140 \pm 0.0007$ \\
0.7 & 0.180 & 0.035 & 0.025 & 1.9 & $0.189 \pm 0.010$ & $0.1804 \pm 0.0005$ \\
\bottomrule
\end{tabular}
\end{table}

\begin{enumerate}[leftmargin=*,itemsep=2pt]
\item \textbf{One damage state} varies by about 0.047 at $f=0.5$. This spread is
what the shaded bands show. Even 200 independent queries would vary by 0.032.
The observed variance is 2.2 times larger, because queries share edges: the
200 queries use only 39 drugs, an average of 5.1 queries per drug, so one lost
drug--protein edge can break several queries at once. The benchmark behaves
like about 90 independent queries. This is a property of the task, not an
error of the model.
\item \textbf{The mean of the ten fixed orders} has a standard error of about
0.013. These means lie 1.4--2.5 points above the formula at all six levels with
$0.1\le f\le0.6$, which is 1.3--2.3 standard errors. At a fixed $f$ the ten
orders are independent, so the standard error is valid. Across levels, however,
the damage states of one order are nested prefixes, so an order that is lucky
at $f=0.3$ tends to stay lucky at $f=0.5$ (their deviations correlate at 0.58).
Six deviations of the same sign are therefore largely one sampling
fluctuation, not six independent ones.
\item \textbf{The expectation}, estimated from 5,000 independent orders, has a
standard error below 0.0007. It matches Eq.~\eqref{eq:r2} to within 0.001.
\end{enumerate}

\subsection{Which query fails?}
\label{sec:blind-perquery}

Given Eq.~\eqref{eq:r2}, the outcome of a query under blind damage depends only
on $m$. No other property of
the intact graph can predict which of two queries with the same $m$ will fail.
We tested this on damage orders 7--9. The formula ranks failures with an AUC of
0.713 at $f=0.3$ and 0.720 at $f=0.5$. A logistic regression on $m$ and four
degree features, trained on orders 0--6, reaches 0.712 at both levels. The expected AUC if outcomes
were drawn from the formula probabilities is 0.72. The redundancy count is thus
a complete predictor under blind damage. Section~\ref{sec:targeted-retrieval}
shows that this changes under attack.
